\section*{第 \ref{ch:g10:heart-lungs-effort} 章——运动中的心脏与肺}

\begin{solution}{exo:g10:heart-lungs-effort:1}
$0.6 \times 14 = \qty{8.4}{L/min}$; $2.2 \times 35 = \qty{77}{L/min}$.
\end{solution}

\begin{solution}{exo:g10:heart-lungs-effort:2}
心脏单侧每分钟泵出的血液体积：
$Q = f \times V_s$. $65 \times 0.075 \approx \qty{4.9}{L/min}$.
\end{solution}

\begin{solution}{exo:g10:heart-lungs-effort:3}
右心房、右心室、肺动脉、肺（肺循环）、
肺静脉、左心房、左心室、主动脉、
器官（体循环）、静脉、右心房。
\end{solution}

\begin{solution}{exo:g10:heart-lungs-effort:4}
潮气量 \qty{0.5}{L}；肺活量 \qty{4.5}{L}；
残气量 \qty{1.2}{L}。
\end{solution}

\begin{solution}{exo:g10:heart-lungs-effort:5}
约 $220 - 16 = 204$ 和 $220 - 60 = 160$ 次/分钟。
\end{solution}

\begin{solution}{exo:g10:heart-lungs-effort:6}
\qty{100}{W}: $120 \times 0.105 \approx \qty{12.6}{L/min}$;
\qty{300}{W}：$198 \times 0.115 \approx \qty{22.8}{L/min}$；两者之比
为 1.8。
\end{solution}

\begin{solution}{exo:g10:heart-lungs-effort:7}
静息：\qty{2.5}{L/min}；运动：$0.04 \times 24 \approx
\qty{1.0}{L/min}$。占比下降了十二倍，绝对流量只
下降 2.6 倍，因为输出量本身已经增加。
\end{solution}

\begin{solution}{exo:g10:heart-lungs-effort:8}
$18 \times 0.14 = \qty{2.5}{L/min}$，约为
训练有素学生 \qty{3.45}{L/min} 的四分之三，
对体型较小者则接近该值。
\end{solution}

\begin{solution}{exo:g10:heart-lungs-effort:9}
大脑即使缺氧和缺葡萄糖数秒也无法耐受，因此
其供应受到优先保障。肠道和肾脏可暂停消化和
滤过一小时而无碍，它们的血流份额让给
肌肉。
\end{solution}

\begin{solution}{exo:g10:heart-lungs-effort:10}
$V_s = 5000/42 \approx \qty{120}{mL}$，而一般约为 \qty{70}{mL}；她
增大的心脏只需更少的搏动即可输出静息时的\omterm{def:g10:heart-lungs-effort:output}{心输出量}。
\end{solution}

\begin{solution}{exo:g10:heart-lungs-effort:11}
大脑发出的运动指令，以及对运动负荷的预判，
会在肌肉收缩前传至心脏中枢和肾上腺。
若跑者只依赖对二氧化碳的感知，则起跑时
\omterm{def:g10:heart-lungs-effort:output}{心输出量}仍处于静息水平，前一分钟靠\omterm{def:g10:cell-metabolism:fermentation}{发酵}供能，
并在供血跟上之前因乳酸积累而被迫减速。
\end{solution}

\begin{solution}{exo:g10:heart-lungs-effort:12}
$195 \times 0.110 \times 0.150 \approx \qty{3.2}{L/min}$ 与
$195 \times 0.160 \times 0.150 \approx \qty{4.7}{L/min}$。每搏输出量
正是训练所增大的量：第二位受试者为
耐力训练者。
\end{solution}

\begin{solution}{exo:g10:heart-lungs-effort:13}
吸入的氧气：$0.21 \times 100 = \qty{21}{L/min}$；消耗量：
\qty{3.5}{L/min}，即 17\%。其余部分被再次呼出——呼出气
中仍含约 16\% 的氧气，这
正是口对口人工呼吸有效的原因。
\end{solution}

\begin{solution}{exo:g10:heart-lungs-effort:14}
\omterm{def:g10:heart-lungs-effort:output}{心输出量}至多约为 $110 \times 0.12 \approx \qty{13}{L/min}$，而不是
24：$\dot V\!\mathrm{O_2}$max 下降近一半，至约
\qty{2}{L/min}。患者可以缓慢步行和骑车，但无法维持
任何高强度运动。
\end{solution}

\begin{solution}{exo:g10:heart-lungs-effort:15}
每升血液的输送量下降四分之一，而氧的提取比例
不变：$\dot V\!\mathrm{O_2}$max 下降 25\%。在随后的
数周内，身体生成更多红\omterm{def:g10:cells-common-unit:cell}{细胞}，使
每升血液携氧量回升至接近 \qty{200}{mL}。
\end{solution}

\begin{solution}{pb:g10:heart-lungs-effort:1}
\textbf{1.} $62 \times 0.072 \approx 4.5$; $105 \times 0.100 = 10.5$;
$140 \times 0.116 \approx 16.2$; $175 \times 0.120 = 21.0$;
$198 \times 0.120 \approx \qty{23.8}{L/min}$.

\textbf{2.} $23.8/4.5 \approx 5.3$：心率贡献 $198/62
\approx 3.2$，每搏输出量贡献 $120/72 \approx 1.7$，二者乘积为 $3.2 \times
1.7 \approx 5.3$.

\textbf{3.} 约 \qty{225}{W} 以上时，每搏输出量保持在
\qty{120}{mL}，仅心率的升高使输出量进一步增加。

\textbf{4.} $220 - 17 = 203$：他的 198 在该经验规则的精度范围之内。

\textbf{5.} $5/4.5 \approx \qty{1.1}{min}$，约 \qty{67}{s}；在
\qty{300}{W}时，$5/23.8 \approx \qty{0.21}{min}$，约 \qty{13}{s}。

\textbf{6.} $0.5 \times 13 = 6.5$; $1.2 \times 18 = 21.6$; $1.8 \times
24 = 43.2$; $2.3 \times 32 = 73.6$; $2.6 \times 44 \approx
\qty{114}{L/min}$.

\textbf{7.} $114/6.5 \approx 18$：潮气量升高到 5.2 倍，
呼吸频率升高到 3.4 倍；体积的贡献更大。

\textbf{8.} $2.6/5.0 = 52\%$，即每次呼吸占肺活量的比例。

\textbf{9.} 不是：通气还有很大的储备，离开
肺的血液仍保持饱和。限制因素是心脏的
每搏输出量，它决定每次搏动能输送多少氧气。

\textbf{10.} $0.7 \times 6.5 \approx \qty{4.6}{L/min}$，$0.7 \times
114 \approx \qty{80}{L/min}$。

\textbf{11.} $4.5 \times 0.052 \approx 0.23$；$10.5 \times 0.090
\approx 0.95$；$16.2 \times 0.120 \approx 1.94$；$21.0 \times 0.145
\approx 3.05$；$23.8 \times 0.155 \approx \qty{3.7}{L/min}$。

\textbf{12.} $3700/68 \approx \qty{54}{mL/min}$（每千克）。摄氧
速率已达最大值，从 225 到 \qty{300}{W} 的增加很小：
这是他的 $\dot V\!\mathrm{O_2}$max，属于健康且经过训练的
学生的水平。

\textbf{13.} $3.7/0.23 \approx 16$：心脏的因子 5.3 乘以
氧提取的因子 $155/52 \approx 3.0$。

\textbf{14.} $3.7 \times 20 = \qty{74}{kJ}$（每分钟），约
\qty{1230}{W}；静息时 $0.23 \times 20/60 \approx \qty{77}{W}$；高于
静息的部分约 \qty{1150}{W}，因此效率为 $300/1150 \approx
26\%$。

\textbf{15.} 在 0 到 \qty{225}{W} 之间，从 \num{0.23} 到
\qty{3.05}{L/min} 大致呈一条直线，随后曲线弯曲：最后的
\qty{75}{W} 只增加 \qty{0.65}{L/min}。上限约为
\qty{3.7}{L/min}。

\textbf{16.} 预期反应：大脑对运动强度的预期
被传递到心血管中枢和肾上腺，在任何运动开始之前提高
心率。

\textbf{17.} 即使在最大通气量和最大\omterm{def:g10:heart-lungs-effort:output}{心输出量}下，\omterm{def:g10:heart-lungs-effort:ventilation}{肺泡}也能使血液充分氧合；
肺并非限制因素。限制因素是心脏每分钟能够泵出的
血量。

\textbf{18.} $20/23.8 \approx 84\%$。肠道和肾脏在静息时占去一半的输出量，
现在只能获得约 \qty{1}{L/min}；它们的动脉已经收缩，
血流已被重新导向
肌肉。

\textbf{19.} $198 \times 0.140 \times 0.155 \approx \qty{4.3}{L/min}$，
约 \qty{63}{mL/min}（每千克）。

\textbf{20.} \omterm{def:g10:heart-lungs-effort:output}{心输出量}乘以 5.3，氧提取乘以 3，
通气乘以 18；训练改变了心脏的每搏输出量，
也随之改变了上限。
\end{solution}
